\documentclass[11pt]{article}
\usepackage[margin=0.75in]{geometry}
\usepackage{booktabs}
\usepackage{array}
\usepackage{amsmath}
\usepackage{hyperref}
\usepackage{enumitem}
\setlist{itemsep=0.25em, topsep=0.35em}
\hypersetup{colorlinks=true, linkcolor=black, urlcolor=blue}
\pagestyle{plain}
\setlength{\parskip}{0.45em}
\setlength{\parindent}{0pt}

\title{\vspace{-1.2em}\textbf{Laidlaw Progress Brief}\\[0.35em]
\large Thermal Extremes and Cardiovascular Admissions in Hong Kong}
\author{Bob Shen \quad|\quad Supervisor: Professor David Bishai\\
\small First-stage update \quad 10 July 2026}
\date{}

\begin{document}
\maketitle
\vspace{-1.5em}

\begin{center}
\textit{Purpose of this note: show what is ready, what is real versus provisional,\\
and what decisions unlock the next stage. No Hospital Authority outcome results yet.}
\end{center}

\section*{1. Where things stand}
I am back in Hong Kong and available full-time for this Laidlaw project. Over the past week I built a reproducible first-stage pipeline, processed public Hong Kong Observatory (HKO) climate data for 2013--2023, reviewed the key Hong Kong temperature--CVD literature, and prepared the data requests we will need from Roro (HA) and Hogan (environment/pollution).

\textbf{Central point:} the environmental and coding infrastructure is substantially ready. Cardiovascular models have so far been tested only on synthetic outcomes. The next scientific milestone is to finalize outcome definitions and obtain the real HA extract, C\&SD denominators, and EPD pollution files.

\section*{2. Scientific framing (recommended)}
\begin{quote}
\textit{Thermal extremes and cardiovascular hospital burden in an aging Hong Kong, 2013--2023.}
\end{quote}
The question is whether officially defined hot nights and cold days are associated with AMI, ischemic stroke, and hemorrhagic stroke admissions among residents aged 35+, and whether the contemporary pattern still looks cold-dominant---as in earlier Hong Kong work---or whether heat metrics (especially hot nights) are now more informative. A cold-dominant result would still be publishable and policy-relevant. We should \textbf{not} frame the project as ``proving heat replaced cold.''

\section*{3. What is real so far: HKO thermal extremes}
Using official HKO definitions at the Observatory Headquarters series:

\begin{center}
\begin{tabular}{@{}lcccc@{}}
\toprule
Period & Hot nights/yr & Very hot days/yr & Cold days/yr & Notes \\
\midrule
2013--2018 (avg) & $\sim$31 & $\sim$30 & $\sim$15.5 & Lower heat extremes \\
2019--2023 (avg) & $\sim$53 & $\sim$48 & $\sim$10.4 & More hot nights/VHDs \\
\bottomrule
\end{tabular}
\end{center}

\vspace{0.3em}
\noindent Definitions: hot night = $T_{\min}\ge 28^\circ$C; very hot day = $T_{\max}\ge 33^\circ$C; cold day = $T_{\min}\le 12^\circ$C.

\noindent Additional verified points:
\begin{itemize}
  \item Peak hot-night year in the core window: \textbf{2021 (61 nights)}.
  \item Peak extremely hot-day year ($T_{\max}\ge 35^\circ$C): \textbf{2022 (15 days)}.
  \item Cold days remain variable (range 1--21/year), not simply disappearing.
  \item Monthly hot nights and very hot days are highly correlated ($r\approx 0.88$), so they should not both enter the same default model.
\end{itemize}
These are \textbf{descriptive climate findings only}. They do not yet say anything about hospital admissions.

\section*{4. What I built (workflow, not findings)}
\begin{itemize}
  \item End-to-end R pipeline: HKO download $\rightarrow$ monthly extremes $\rightarrow$ panel merge $\rightarrow$ NB/Poisson templates $\rightarrow$ diagnostics.
  \item Planned primary model: Negative Binomial with offset $\log(\mathrm{population}\times\mathrm{days})$, month effects, time trend, COVID phases; primary heat metric = hot nights; co-primary = cold days.
  \item Synthetic HA-like panel used only to test the code. \textbf{Synthetic coefficients are not study results.}
  \item Draft data-request memos for Roro (HA) and Hogan (climate/pollution).
\end{itemize}

\section*{5. Still missing / recommended defaults}
\begin{itemize}
  \item \textbf{Missing:} HA aggregate extract (Roro); real C\&SD age--sex denominators; real EPD general-station pollution files (Hogan).
  \item \textbf{Defaults I'd like your advice on:} core window 2013--2023 (2024 optional); inpatient principal-diagnosis AMI/ischemic/hemorrhagic stroke; separate ages 65--69 and 70--74; ED separate; ozone staged (not automatic confounder); historical comparison with Goggins-era work kept qualitative first.
\end{itemize}

\section*{6. Ask}
A 30-minute meeting with you (Hogan/Roro welcome) to lock outcome definition, study window, and data-access route. I can move full-time once those are set.

\vspace{0.5em}
\noindent\textit{Repo:} \url{https://github.com/bobshenruililin/Laidlaw-Heat-Project}

\end{document}
